\section{Feature construction}
\begin{table}[htbp]
\centering\small
\caption{Household representations fitted on calibration purchases.}
\label{tab:features}
\begin{tabularx}{\textwidth}{p{28mm}X}
\toprule
Representation & Definitions and preprocessing\\
\midrule
RFM (3 columns) & Recency is $501-\max(\mathrm{DAY})$; frequency is distinct basket count; monetary value is positive recorded sales summed over calibration. Apply $\log(1+x)$, then standardize each column.\\
Behavior (8 columns) & RFM plus mean basket sales, mean basket line count, distinct purchased products, distinct departments and proportion of lines carrying any negative retail/coupon discount field. Log-transform all except the discount proportion, then standardize.\\
Product (17 columns) & Behavior plus sales shares in the eight largest calibration departments and an Other category. Take square roots, center each share coordinate, and scale the composition block to the same total variance as the eight-column Behavior block.\\
\bottomrule
\end{tabularx}
\end{table}
The eight largest departments are selected from calibration sales, not the later period. The Other group retains the remaining sales so the composition sums to one. Square-root shares provide Hellinger-style coordinates: absent departments remain finite without adding arbitrary pseudocounts. Centering does not change pairwise distances; the block scale explicitly controls how much department composition contributes relative to behavior. Equal block variance is a modeling choice; sensitivity to this weight is examined later.

Frequency measures distinct baskets rather than transaction-line count. Mean basket lines counts recorded positive purchase lines rather than distinct items or units. Product breadth and department breadth measure diversity over the whole calibration period, so they remain related to shopping intensity. Mean basket value is also algebraically connected to total sales and basket frequency. Standardization does not remove those relationships, so the enriched representation may still be dominated by purchase intensity.

Retail and coupon discount columns contain negative values in this dataset. They are not subtracted from recorded sales; only their sign is used, as a recorded-discount indicator. Positive-purchase filtering excludes returns and zero-price transactions, so features describe recorded buying activity.

\section{Reproducible analytical workflow}
The executable notebook embeds the analysis implementation, mirrors it in the project source module, and regenerates all reported tables and figures. The workflow first verifies the two required input files, loads transactions and products, checks product-key uniqueness, and performs a validated many-to-one left join. The joined row count must match the transaction count. It records unmatched products, missing values, nonpositive quantities and sales, duplicate source rows, and basket consistency checks.

Exact duplicate source rows are retained and counted, because repeated item scans cannot be distinguished from erroneous duplication without line identifiers. Join duplicates are prevented by key validation and row-count checks. Missing departments, including blank text, receive an explicit Unknown category. Other exclusion and audit counts are reported in the following section.

The same cohort and derived matrices feed K-means, Ward hierarchical clustering and DBSCAN. The K-means and Ward grid uses two through eight clusters. K-means uses 20 initializations for the main fit with random seed 42. DBSCAN uses a minimum of ten samples and neighborhood radii chosen from the 50th, 70th, 85th and 95th percentiles of tenth-neighbor distances (including the observation itself among the ten points) in each space.

A representative K-means fit is chosen separately within each space by highest within-space silhouette among partitions with every group containing at least five percent of the cohort. The rule selects a model within a space, not across spaces. Repeated-initialization tests use 20 one-initialization fits. Subsample tests use 20 paired independent samples, each containing 80 percent of households, and compare assignments only for shared households. Preprocessing is held fixed in these checks.

Later positive purchase sales and distinct baskets are summarized for every calibration household. Households without later recorded purchases receive zero, retaining the cohort rather than conditioning on return shoppers. Mean-spend uncertainty uses 1,000 within-cluster bootstrap resamples, giving percentile intervals conditional on the fitted groups. Between-group explained variance ($\eta^2$) depends on group count, so it is read alongside the paired fixed-count comparison.
